\clearpage
\section{Model prototypu}

    \begin{figure}[!h]
        \centering \includegraphics[width=1\linewidth]{Model/model.pdf}
        \caption{Model prototypu}
    \end{figure}

    \paragraph{}
        W strukturze projektowanego BSP zdecydowano się na zastosowanie układu centropłata, 
        Wybór ten podyktowany jest przede wszystkim pozwoli osiągnąć optymalne osiągi aerodynamiczne oraz wysokią manewrowości
        oraz wyposazony w chowane podwozie o budowie,  
        Konfiguracja ta wiąże się z ograniczeniem przestrzeni użytecznej wewnątrz kadłuba,
        Omawiany dron nie jest przeznaczony do przewozu wielkogabarytowego ładunku płatnego, 
        lecz do obserwacji, 
        Wyposażenie stanowi jedynie zintegrowana głowica optoelektroniczna,
    
    \paragraph{}
        Niemniej jednak, mając na uwadze konieczność optymalizacji czasu trwania procesu projektowego oraz dążenie do minimalizacji ryzyka błędu, 
        zdecydowano się na implementację usterzenia klasycznego, 
    
    \paragraph{}
        W projekcie zastosowano układ ze śmigłem pchającym z tyłu kadłuba, 
        Decyzja ta wynika z konieczności zapewnienia niezakłóconego pola widzenia dla głowicy optoelektronicznej, 
        Przeniesienie napędu na ogon uwalnia przestrzeń na dziobie drona, 
        eliminując ryzyko wchodzenia łopat śmigła w kadr i powstawania zakłóceń obrazu, 
        Dodatkowo chroni to delikatną optykę przed zanieczyszczeniami i uszkodzeniami podczas lądowania,